\cgeentry{cge_gdx}{GAMS data exchange (GDX) files: write, read and list}

\cgesec{Description}

\cmd{cge_gdx} (from version 0.10.0) exchanges data with GAMS through GDX
files, the binary files of GAMS's \cmd{execute_unload} and \cmd{\$gdxin}.
\cmd{cge_gdx write} writes the calibrated model and its simulations: the SAM,
the roles of its accounts, the calibrated parameters and the levels of the
variables. \cmd{cge_gdx read} reads any symbol of any GDX file into the data in
memory, a parameter, a set or the levels and marginals of a variable, and
\cmd{cge_gdx list} lists the symbols of a file. \manref{cge_sam} reads a SAM
from a GDX file. GAMS writes the files and its tool \cmd{gdxdump} reads them,
so GAMS must be installed; the numbers pass to the last digit.

\cgesec{Quick start}

The model and a simulation as a GDX file
\begin{stcode}
cge_gdx write using model.gdx, sims(notariff) gams(C:/GAMS/54/gams.exe)
\end{stcode}
Its symbols, and the levels of the variables read back
\begin{stcode}
cge_gdx list using model.gdx, gams(C:/GAMS/54/gams.exe)
cge_gdx read using model.gdx, symbol(LEVEL) clear gams(C:/GAMS/54/gams.exe)
\end{stcode}
A SAM from the GDX file of another GAMS model
\begin{stcode}
cge_sam using other.gdx, symbol(SAM) clear gams(C:/GAMS/54/gams.exe)
\end{stcode}
The solution of the GAMS program of \manref{cge_gams}, written by GAMS with
\cmd{gdx=sol.gdx}
\begin{stcode}
cge_gdx read using gams_notariff/sol.gdx, symbol(QA) clear gams(C:/GAMS/54/gams.exe)
\end{stcode}

\cgesec{Syntax}

\begin{stsyntax}
\cmdb{cge_gdx write} \cmd{using} \argu{file}\cmd{.gdx} [\texttt{,} \optarg{}{sims}{names} \opt{}{replace} \opt{}{norun} \optarg{}{gams}{path}]\\
\cmdb{cge_gdx read} \cmd{using} \argu{file}\cmd{.gdx} \texttt{,} \optarg{sym}{bol}{name} \opt{}{clear} [\opt{}{norun} \optarg{}{gams}{path}]\\
\cmdb{cge_gdx list} \cmd{using} \argu{file}\cmd{.gdx} [\texttt{,} \opt{}{clear} \opt{}{norun} \optarg{}{gams}{path}]
\end{stsyntax}

\begin{opttable}
\optgroup{Main}
\optarg{}{sims}{names} & the simulations whose levels are written, with the base; \cmd{_all} for every one\\
\optarg{sym}{bol}{name} & the symbol to read\\
\opt{}{clear} & replace the data in memory\\
\opt{}{replace} & replace an existing GDX file\\
\optgroup{GAMS}
\opt{}{norun} & write the script without running it; read the output of an earlier run\\
\optarg{}{gams}{path} & the GAMS executable, with \cmd{gdxdump} in its folder\\
\end{opttable}

\cgesec{Options}

\begin{optdesc}
\item[\optarg{}{sims}{names}] names the stored simulations whose levels are
written, with the base; \cmd{_all} writes every one, and by default only the
base is written.

\item[\optarg{sym}{bol}{name}] names the symbol of the file to read.

\item[\opt{}{clear}] replaces the data in memory, which \cmd{cge_gdx read}
needs; with \cmd{cge_gdx list}, the list is left as a dataset.

\item[\opt{}{replace}] allows an existing GDX file to be replaced.

\item[\opt{}{norun}] writes the script without running it; a command that reads
then reads what an earlier run of its script wrote.

\item[\optarg{}{gams}{path}] names the GAMS executable, such as
\cmd{gams(C:/GAMS/54/gams.exe)}; \cmd{gdxdump} is taken from its folder. By
default \cmd{gams} and \cmd{gdxdump} must be on the system path.
\end{optdesc}

\cgesec{Remarks and examples}

\cgesubsec{What is written}

\cmd{cge_gdx write} writes a GAMS program, \argu{file}\cmd{_write.gms},
which holds every number to the last digit and ends with
\cmd{execute_unload}, and GAMS runs it. The GDX file holds:
\begin{cgetable}{>{\ttfamily}l p{0.66\linewidth}}
\toprule
\normalfont Symbol & Content\\
\midrule\endhead
\bottomrule\endlastfoot
AC & the accounts of the SAM\\
A, C, F, H, E & the activities, commodities, factors, households and enterprises\\
SAM(AC, AC) & the SAM, after the layout rewrites of calibration\\
CAL(*, *, *) & the calibrated parameters: name, element, element\\
sim & the base and the simulations written\\
LEVEL(sim, *, *, *) & the levels of the variables: simulation, variable, element, element\\
\end{cgetable}
An index that a parameter or variable does not have is \cmd{-}, so that
\cmd{CAL('tm', 'com\_man', '-')} is the tariff on manufactures and
\cmd{LEVEL('base', 'EXR', '-', '-')} the base exchange rate. GAMS stores no
zeros.

\cgesubsec{How the files are read}

\cmd{gdxdump} writes a symbol as a CSV file, with its numbers as hexadecimal
floats, which \cmd{cge_gdx read} reads exactly. The data in memory then hold
one observation per record, with a string variable for each index, named
after its domain (\cmd{dim1}, \cmd{dim2}, \dots\ for an index without one), and
the values: \cmd{value} for a parameter or scalar; \cmd{level},
\cmd{marginal}, \cmd{lower}, \cmd{upper} and \cmd{scale} for a variable or an
equation; \cmd{text} for a set. GAMS's \cmd{EPS} is read as 0, and the
infinities and the special values \cmd{NA} and \cmd{UNDF} as missing.

\cgesubsec{Scripts and batch mode}

Each command writes a script in the folder of the file that runs GAMS or
\cmd{gdxdump} there, and runs it with \cmd{shell}. Stata for Windows cannot
start other programs in batch mode: there, and with \opt{}{norun}, the script is
written but not run. Run it, then give the command again with \opt{}{norun}: it
reads what the script wrote. The examples of this entry were made in that
way, the scripts run between two runs of the do-file.

The SAM is mapped, calibrated and solved:
\stlog{gdx_setup}

\cgesubsec{Example 1: writing a GDX file}

\stlog{gdx_write}

\cgesubsec{Example 2: the symbols of a file}

\stlog{gdx_list}

\cgesubsec{Example 3: a SAM from a GDX file}

\manref{cge_sam} reads a parameter with two indexes, \cmd{SAM} by default, as
long data:
\stlog{gdx_sam}

\cgesubsec{Example 4: reading a symbol}

The levels written, read back to the last digit:
\stlog{gdx_read}

\cgesec{Stored results}

\cmd{cge_gdx write} stores the following in \cmd{r()}:
\begin{storedres}
\resgroup{Scalars}
\res{r(N_sims)} & number of simulations written\\
\res{r(N_levels)} & number of levels written\\
\res{r(N_parameters)} & number of parameter values written\\
\resgroup{Macros}
\res{r(sims)} & the base and the simulations\\
\end{storedres}
\cmd{cge_gdx read} stores:
\begin{storedres}
\resgroup{Scalars}
\res{r(N)} & number of records\\
\res{r(dim)} & number of indexes\\
\resgroup{Macros}
\res{r(kind)} & \cmd{parameter}, \cmd{variable} or \cmd{set}\\
\res{r(indexes)} & the names of the index variables\\
\res{r(symbol)} & the symbol\\
\end{storedres}
\cmd{cge_gdx list} stores \cmd{r(N_symbols)} and \cmd{r(symbols)}.

\cgesec{Methods and formulas}

A number $x$ is written in the GAMS program with the fewest digits, of 15, 16
or 17, that give back the same double, so GAMS reads $x$ itself. \cmd{gdxdump},
with \cmd{dFormat=hexponential}, writes $x = \pm\,0\text{x}h_0.h_1\dots h_k\text{p}e$,
which is read as $\pm\,m \cdot 2^{e - 4k}$, with $m$ the integer of the
hexadecimal digits $h_0 h_1 \dots h_k$, exact in double precision; the round
trip from Stata to GAMS and back therefore gives every number unchanged, as
the check in \texttt{tools/gams/run\_gdx\_check.py} confirms.

\cgesec{Also see}

\begin{alsosee}
\seeentry{cge_gams}{Write the model as a GAMS program, and solve it in GAMS}
\seeentry{cge_sam}{Read and check a social accounting matrix}
\seeentry{cge_parameters}{The calibrated parameters as a dataset}
\seeentry{cge_results}{Report simulation results}
\end{alsosee}
